\documentclass{article}
\PassOptionsToPackage{numbers,sort&compress}{natbib}
\usepackage[dblblindworkshop]{neurips_2026}
\workshoptitle{Foundation Models for Temporal Systems: From Forecasting to World Modeling}
\usepackage[utf8]{inputenc}
\usepackage[T1]{fontenc}
\usepackage{amsmath,amssymb,graphicx,booktabs,longtable,array,calc}
\usepackage{microtype,url}
\usepackage[hidelinks,unicode]{hyperref}
\providecommand{\tightlist}{\setlength{\itemsep}{0pt}\setlength{\parskip}{0pt}}
\providecommand{\pandocbounded}[1]{#1}
\setlength{\emergencystretch}{1em}
\title{Hidden-Feature Matching and Left-Right Movement in Gait Representations}
\author{Anonymous Author(s)}
\hypersetup{pdftitle={Hidden-Feature Matching and Left-Right Movement in Gait Representations},pdfauthor={},pdfcreator={LaTeX}}
\begin{document}
\maketitle
\begin{abstract}
Temporal representations should retain information that helps describe
an evolving system, alongside the features they learn to predict. We
examine this requirement in gait using a small skeleton Joint-Embedding
Predictive Architecture (JEPA). The model predicts hidden skeleton
features; a separate regression estimates a signed contrast between left
and right movement speeds from its frozen encoder. On 625 clips from 93
held-out-by-source videos, trained predictors more consistently match
hidden features from the correct clip than from another video. Yet with
an expanded feature summary, mean source-balanced \(R^2\) for the
movement estimate falls from 0.22 at initialization to 0.10--0.11 after
training. Motion-weighted and connected-region masks do not establish an
improvement over their own random references. This focused negative
finding connects latent prediction to evaluation of observable movement:
success on hidden-feature matching did not improve the specified
movement estimate under the tested training and readout procedure.
Future-only prediction remains a separate test.
\end{abstract}
\section{Introduction}\label{introduction}

Walking is an evolving system: coupled body landmarks change position
over time, and movement depends on their trajectories. Rhythm, phase and
coordination make gait an appealing setting for temporal modeling,
although video-derived coordinates describe only part of the physical
process. They omit forces and contact and are affected by camera view
and pose-estimation error. Temporal world modeling aims to represent
such a system well enough to predict its observable behavior. A
practical evaluation question is whether a representation trained
through prediction preserves information needed for a specified movement
measurement.

We study that question with S-JEPA \citep{sjepa}, using a left--right
contrast of median coordinate speeds. The contrast is clinically
motivated by work on gait asymmetry \citep{stroke}, but is not a
validated clinical index. Our central null is that JEPA training
provides no improvement over the same encoder at initialization when a
frozen readout estimates this contrast on held-out source videos. We
investigate input preparation, hidden-feature correspondence, masking,
and movement readout as distinct parts of that question. The current
expanded-readout and mask comparisons are exploratory analyses on a
repeatedly examined cohort; they do not retrospectively become
registered tests.

Predictive skeleton methods and temporal foundation models share an
evaluation problem: the training target is a learned feature, whose
successful prediction need not preserve every observable of interest.
Our contribution is a controlled, domain-specific test of this issue. It
is relevant to evaluation and reliability, grounded in temporal data and
a temporal encoder, without establishing foundation-model scale, broad
transfer, persistent state or a gait simulator.

\section{Study design}\label{study-design}

\textbf{Data and observable.} Our selected GAVD subset \citep{gavd}
contains 666 annotations, 642 matching pose archives and, after quality
checks, 625 clips from 93 source videos. Each archive has shape
\(T_i\times33\times4\): three estimated coordinates plus landmark
visibility, with original frame numbers and frame rate. Coordinate units
are image normalized with inferred depth; they are not calibrated
physical speeds.

The movement path uses observed coordinates, pelvis centering and a
body-width scale, without gap filling or resizing. For five
pairs---shoulders, knees, ankles, heels and foot tips---we calculate
median speeds \(m_{L,k},m_{R,k}\) using original timestamp differences.
Both sides must be observed at both ends of at least eight common
transitions for every pair. The signed target is

\[y=\frac{1}{5}\sum_{k=1}^{5}\frac{m_{L,k}-m_{R,k}}{m_{L,k}+m_{R,k}+10^{-8}}.\]

Anatomical reflection exchanges left/right landmarks and validity flags
and flips the centered horizontal coordinate, giving \(y(Mx)=-y(x)\).
Hips define the pelvis and are excluded from the contrast because
centering equalizes their speed magnitudes. The statistic does not
measure phase, stride timing or loading; reversing a trajectory together
with its sampling intervals preserves its speed magnitudes.

\textbf{Encoder input.} A landmark is observed when visibility is at
least 0.45 and all coordinates are finite. A separate encoder path fills
interior gaps of at most four samples, treating filled positions as
valid. After pelvis centering and body-width normalization, coordinates
and validity are linearly resized to 64 positions by sequence index.
Original time intervals are discarded. Resized validity must reach
0.999; invalid coordinates are zero. Each joint's four-position patch
projects 12 coordinates to 96 features, giving a \(16\times33\) grid. A
patch is valid only if all four positions are valid. Anatomical
identities and prepared positions survive this operation; the original
trajectory, duration and cadence need not.

\begin{figure}[htbp]
\centering
\includegraphics[width=\linewidth]{assets/v1/figures/training_pipeline_compact.pdf}
\caption{Two measurement paths and the completed JEPA training grid. Poses and mask cells are schematic. The movement target reaches only the separate regression evaluation. The online encoder has visible context on either side of a hidden position; the teacher and full-view regularizer access complete prepared clips.}
\end{figure}

\textbf{Training.} The completed grid uses 96 features,
encoder/predictor depths 4/2, four attention heads, batch size 20 and
1,200 AdamW updates, with learning rate 0.001 and weight decay 0.05. The
online encoder zeroes hidden patch content before adding joint and
position embeddings; the predictor substitutes learned mask tokens at
target positions. Cross-entropy matches centered teacher feature
distributions at hidden valid positions, averaged within clips and then
across clips. The teacher has no gradients and follows the online
weights by exponential moving average (momentum 0.999). A full-view
VICReg term \citep{vicreg}, weighted 0.05, pools twelve gait joints
across two geometric views. Thus the total objective also accesses
complete clips. Training uses CUDA BF16 and evaluation FP32.

We compare three motion conditions and two region conditions across five
folds and five seeds: 125 encoders. The motion conditions use random
sampling, a MAMP-style sampler \citep{mamp}, and a robust motion
mixture. Their target count is half the batch's smallest valid
twelve-joint gait pool, rounded down, sampled across all 33 joints.
Connected regions hide six connected landmarks over eight of sixteen
blocks, with a separate count-matched random reference. Starting
weights, source draws, geometric views and update counts are matched
within each comparison. No reflection augmentation or reflection loss
appears in this grid.

\textbf{Held-out readout.} Five outer folds keep each video's clips and
derived views together. Fresh encoders fit only outer-training videos
with source-balanced sampling. All overlapping windows, mirrors and
augmentations inherit the source split. The encoder is frozen while
ridge regression, a linear model with a coefficient penalty, estimates
\(y\). Three inner source folds select the penalty using only
outer-training labels; imputation and scaling are fitted within each
inner training set and then refitted with the regression on all
outer-training sources. Inner-validation clips may have contributed
without labels to outer-training encoder fitting.

The compact readout uses left-minus-right and left-plus-right mean
features from five pairs (960 inputs). The expanded readout adds feature
standard deviations, absolute changes between adjacent valid blocks, and
ten support fractions (2,890 inputs). Both initial and trained encoders
use this procedure. We pool outer-test predictions within each seed,
weighting each video equally, then average the five \(R^2\) values.
Paired 95\% source-bootstrap intervals resample videos with their clips
and seed predictions together (2,000 resamples). They condition on the
fitted models and existing splits. Repeated seeds and mask draws do not
add independent sources.

\section{Findings}\label{findings}

\textbf{Movement readout falls after training.} The expanded initial
representation reaches mean \(R^2=0.223\), compared with
\(0.101\)--\(0.114\) for trained teachers (Figure 2). All five
conditions decrease in every seed, with greater mean absolute error.
Paired trained-minus-initial intervals are below zero in every
condition. For the random motion mask, the difference is \(-0.109\) with
interval \([-0.170,-0.041]\). Final online encoders also score below
their matched initial encoders.

The compact initial readout scores 0.071; expanding it increases \(R^2\)
by 0.152. The added components mix variation, feature changes and
observation support, so the gain cannot be assigned specifically to
temporal order. The initial control retains the pretrained pose
detector, input preparation, anatomical identities and supervised
regression. A 264-value direct-pose summary reaches \(R^2=0.035\); this
limited baseline does not bound stronger coordinate readouts.

\begin{figure}[htbp]
\centering
\includegraphics[width=\linewidth]{assets/v1/figures/learning_results.pdf}
\caption{Different outcomes from the same fitted models. A: frozen movement readout; points are five seed scores, and short vertical marks their mean. B: correct-clip feature error is lower in 375 of 375 trained diagnostic checks, versus 33 of 75 distinct initial checks. These checks reuse models and videos.}
\end{figure}

\textbf{Correct-clip features become easier to match.} Each diagnostic
compares source-weighted mean squared feature error against the correct
clip and valid targets from another source, holding the evaluated clips
and hidden positions fixed. All 375 trained checks favor correct
targets; 33 of 75 initial checks do so. The denominator is models times
evaluation masks, not successful predictions of 375 independent
movements. This preference could involve pose, viewpoint or missingness
as well as motion. Teachers and target distributions change during
training, so raw training losses cannot rank representation quality
across conditions.

\textbf{Changing mask geometry does not establish better readout.} The
inspection draws confirm that motion samplers select higher-motion
targets and that connected regions remove immediate temporal brackets.
Motion masks hide approximately 17\% of valid all-landmark tokens;
regions hide 9.9\%. Their readout comparisons therefore use distinct
random references. MAMP-style, motion-mixture and region differences are
respectively \(-0.0008\), \(+0.0005\) and \(-0.009\), with intervals
\([-0.020,0.015]\), \([-0.015,0.015]\) and \([-0.033,0.018]\). Intervals
spanning zero leave benefits uncertain, rather than establishing
equivalence.

\textbf{Preparing input changes the measurement.} A saved comparison
recalculates the target after input preparation: 623 clips from 92
sources have finite values on both paths, with source-weighted sign
agreement 70.4\% and calculation-agreement \(R^2=0.218\). The two paths
disagree, but this calculation neither sets a model-performance ceiling
nor identifies an individual preprocessing cause.

\section{Interpretation and next
experiment}\label{interpretation-and-next-experiment}

The two evaluations answer different questions. JEPA training improves
preference for the correct clip's hidden features, while a fixed
regression procedure estimates its left--right speed contrast less
accurately. Poorer readout performance does not imply complete
destruction of movement information. In particular, 49/125 teacher fits
and 96/125 online fits select the largest ridge penalty, 10,000. Wider
training-only selection, support-only controls and summaries of matched
dimension would help separate feature geometry from readout limitations.
Person identities, reuploads, pose errors and recording cues remain
possible explanations beyond what source holdout controls.

A next experiment would predict future coordinates from a past-only
prefix at common valid endpoints and physical-time horizons, against
last-position, constant-velocity, direct-pose and initial-encoder
controls. Any decoder and settings must be fitted on training sources
only. Correctly paired training futures should be compared with
mismatched futures. Future coordinates or visibility must not affect
prefix preparation, normalization, mask selection or forecasts. Existing
past-feature readout and future-feature decoding demonstrations have
synthetic training evidence; a real-data horizon census establishes
eligibility only. Decoding observed future-teacher features diagnoses
the decoder; decoding predictor features produced from the past tests
forecasting. Multiple horizons from one prefix are not recursive rollout
evidence.

This study supports evaluating latent prediction alongside access to
specified observable behavior. The evidence remains conditional on one
cohort, JEPA recipe and readout, with no participant-independent or
distribution-shift evaluation. Forecasting and simulation would demand
stronger tests of future behavior and consistency than this whole-clip
contrast supplies.

\section{Related work}\label{related-work}

S-JEPA evaluates hidden-feature prediction through skeleton action
recognition \citep{sjepa}, whereas MAMP reconstructs motion and uses
motion-aware masking \citep{mamp}. Connected-body masking in SLiM aims
to reduce local shortcuts \citep{slim}. Our controlled masks borrow
these ideas while holding the JEPA objective fixed; they do not
reproduce those complete methods. V-JEPA 2 couples video representations
with anticipation and action-conditioned planning \citep{vjepa2}, and
seq-JEPA separates transformation-sensitive from invariant features
\citep{seqjepa}. Together these works motivate evaluation beyond latent
loss. Our narrower contribution compares matched initial and trained
features on an observed gait quantity, without claiming their
forecasting or world-modeling capabilities.

\clearpage
\bibliographystyle{unsrtnat}
\bibliography{assets/v1/references}
\clearpage
\appendix
\section{Measurement and implementation
details}\label{measurement-and-implementation-details}

The archive coordinates are \(x=(x_0+uW_c)/W\), \(y=(y_0+vH_c)/H\), and
\(z=dW_c/W\), where crop origin and size are \((x_0,y_0),W_c,H_c\),
frame size is \(W,H\), and \(u,v,d\) are pose estimates. Unequal
image-axis scaling and inferred depth remain after body normalization.
The scale is the median of valid shoulder and hip widths. The movement
path requires observed hips; the encoder uses the median valid pelvis
midpoint for missing hips (zero if none is available). No boundary
extrapolation fills gaps.

The teacher receives the unaugmented prepared clip; the masked online
path receives a geometric view. Two full online views use rotations up
to 8 degrees in the horizontal-depth plane and translations up to 0.03
in image-plane coordinates. Target/predictor temperatures are 0.06/0.10;
the teacher center uses momentum 0.9. Gradient norm is clipped at 1. The
twelve-joint regularizer pool is shoulders, hips, knees, ankles, heels
and foot tips. Its projector is a separate training component. It
combines invariance, variance and covariance penalties with internal
weights 25, 25 and 1, then overall weight 0.05.

The expanded summary uses common left/right validity within each pair;
increments require consecutive common valid blocks. Each pair
contributes differences and sums of means, standard deviations and mean
absolute feature increments (six 96-vectors), plus two support
fractions. Thus \(5\times6\times96+10=2,890\). Missing support gives
zero summary values. Temporal positions in this summary are prepared
block indices, not original timestamps. Ridge penalties are selected
from \(0.01,0.1,1,10,100,1000,10000\).

Source-balanced \(R^2\) uses weights \(w_{vi}=1/(Vn_v)\) and weighted
test mean \(\bar y_w\):

\[R_w^2=1-\frac{\sum_{vi}w_{vi}(y_{vi}-\hat y_{vi})^2}{\sum_{vi}w_{vi}(y_{vi}-\bar y_w)^2}.\]

The test mean defines this metric's denominator, not a deployable fitted
baseline. Outer test folds contain 19, 19, 19, 18 and 18 videos and 189,
182, 72, 77 and 105 clips. Source IDs and split membership are checked;
no persistent participant identities or complete cross-upload content
comparison establish independence across different IDs. Annotation
retrieval/version records are incomplete.

\section{Supporting results and
reproducibility}\label{supporting-results-and-reproducibility}

{\def\LTcaptype{none} % do not increment counter
\begin{longtable}[]{@{}
  >{\raggedright\arraybackslash}p{(\linewidth - 6\tabcolsep) * \real{0.2500}}
  >{\raggedleft\arraybackslash}p{(\linewidth - 6\tabcolsep) * \real{0.2500}}
  >{\raggedleft\arraybackslash}p{(\linewidth - 6\tabcolsep) * \real{0.2500}}
  >{\raggedright\arraybackslash}p{(\linewidth - 6\tabcolsep) * \real{0.2500}}@{}}
\toprule\noalign{}
\begin{minipage}[b]{\linewidth}\raggedright
Training condition
\end{minipage} & \begin{minipage}[b]{\linewidth}\raggedleft
Mean \(R^2\)
\end{minipage} & \begin{minipage}[b]{\linewidth}\raggedleft
Trained minus initial
\end{minipage} & \begin{minipage}[b]{\linewidth}\raggedright
95\% paired source interval
\end{minipage} \\
\midrule\noalign{}
\endhead
\bottomrule\noalign{}
\endlastfoot
Random, motion budget & 0.114 & -0.109 & {[}-0.170, -0.041{]} \\
MAMP-style & 0.113 & -0.110 & {[}-0.165, -0.051{]} \\
Motion mixture & 0.114 & -0.108 & {[}-0.171, -0.042{]} \\
Random, region budget & 0.109 & -0.113 & {[}-0.166, -0.060{]} \\
Connected region & 0.101 & -0.122 & {[}-0.183, -0.055{]} \\
\end{longtable}
}

These exploratory percentile intervals use bootstrap seed 812, retain
all five fitted seeds and have no familywise or development-selection
adjustment. The numerical supplement preserves aggregate seed scores and
recorded intervals. The executed report identifies 125,000 prediction
rows and 200 pooled score rows; a saved audit records their
recomputation. Raw predictions and checkpoints are unavailable in the
revision environment, so this revision checks aggregate consistency
without claiming to repeat that bootstrap or inference.

A separate reflection-augmentation comparison has real-data summary
evidence: augmentation reduces the specified teacher-token reflection
discrepancy by \(\Delta q=-0.008\) \([-0.010,-0.007]\), while its
readout difference is \(\Delta R^2=0.004\) \([-0.006,0.013]\). This uses
its own recipe, four inner readout folds and compact summary; it is not
pooled with the motion grid. An explicit reflection loss has synthetic
evidence only. Likewise, a readout forced to take the form
\(g(x)=[h(x)-h(Mx)]/2\) reverses sign algebraically without establishing
learned movement information.

GAVD annotations describe the cohort, not affected side or treatment
response. Uncalibrated coordinates and absent clinical reference
measurements preclude a clinical interpretation of \(y\). Public
annotation access does not authorize redistribution of separately hosted
videos or derived poses; neither is included in the supplement.

\textbf{AI assistance.} An AI assistant supported author-side evidence
reconciliation, literature checking, manuscript revision, figure
programming and adversarial checks. No new model training was performed
for this revision. Responsibility for the submitted text, figures,
claims and references remains with the authors.

\end{document}
